\section{Introduction}
\label{sec:introduction}

FinOps practice becomes materially harder once cost governance spans multiple cloud providers. AWS, Azure, and GCP each expose different billing exports, discount semantics, commitment accounting behaviors, and ownership signals \cite{aws_line_items,azure_amortized_cost,gcp_standard_export}. A team that looks inexpensive in one raw export can appear expensive in another simply because reservation amortization, savings-plan effective cost, or credit handling is encoded differently. The first systems challenge in multi-cloud optimization is therefore not recommendation ranking or dashboard design; it is establishing a cost model that makes cross-cloud comparison economically and operationally defensible.

This problem is especially acute for post-billing decision support. Raw billing data mixes consumed workload cost with idle commitment waste, hides responsibility behind incomplete tags or shared subscriptions, and often collapses distinct savings notions into one unsafe number. In the current benchmark state of this repository, the unified billing mart contains 18{,}992 normalized rows across AWS, Azure, and GCP, with \$52{,}021.31 in net effective cost (NEC) and \$29{,}874.07 of unattributed NEC. These values are useful precisely because they are derived from an explicit canonical model rather than from provider-local billing conventions.

Most industry tooling provides descriptive visibility first: spend breakdowns, trend charts, and vendor-specific optimization hints. Those views are useful, but they are not enough for auditable multi-cloud decision making. A systems paper must instead answer four harder questions: how heterogeneous billing exports are normalized into one data contract, how commitment-aware effective cost is computed consistently, how shared and unattributed spend become accountable optimization scope, and how recommendations, forecasts, and validation outputs are generated transparently from that canonical substrate.

This paper presents a reproducible post-billing \emph{multi-cloud FinOps decision engine} rather than a dashboard-only cost viewer. The system ingests provider billing exports, computes cloud-specific NEC semantics, materializes a Unified Cost Allocation Schema (CAS), allocates shared spend into accountable scope, detects optimization and governance issues, ranks recommendations with auditable scoring, forecasts month-end NEC, and persists lifecycle-aware validation outputs. All stages are implemented as reproducible repository artifacts over a synthetic multi-cloud benchmark, allowing the paper to make concrete design and evaluation claims without depending on proprietary enterprise billing data.

\subsection*{Contributions}
The paper makes four primary contributions:
\begin{itemize}
    \item It defines and implements a canonical Unified Cost Allocation Schema (CAS) that normalizes heterogeneous AWS, Azure, and GCP billing semantics into one downstream data contract for attribution, NEC analysis, and decision generation.
    \item It introduces a commitment-aware NEC modeling layer that separates consumed workload cost from commitment waste and avoids the failure modes of naive billing fields.
    \item It presents an auditable decision pipeline that connects waste detection, structured reasoning over billing facts, explicit signal-to-action mapping, deterministic recommendation scoring, and lightweight month-end forecasting.
    \item It demonstrates the system on a reproducible synthetic benchmark, yielding 95 materialized recommendations and forecast backtests with \$7.86 mean absolute error, 2.98\% mean absolute percentage error, and a 95.59\% within-10-percent rate.
\end{itemize}

Two motivating examples clarify why this framing matters. First, RI-covered or Savings Plan-covered AWS usage can appear misleadingly cheap or free in raw fields, even though the enterprise has already incurred commitment cost. Second, tag gaps and shared-cost rows can block accountable optimization even when total spend is visible; the current benchmark's \$29{,}874.07 in unattributed NEC shows that visibility without attribution is not sufficient for action.

The remainder of the paper is organized as follows. Section~\ref{sec:cas} defines the Unified Cost Allocation Schema and its cross-cloud mapping decisions. Section~\ref{sec:allocation} explains shared-cost allocation and accountable cost propagation. Section~\ref{sec:nec} formalizes NEC modeling and validation. Section~\ref{sec:intelligence} presents the decision-intelligence pipeline, including waste detection, recommendation scoring, forecasting, and lifecycle validation. Section~\ref{sec:eval} summarizes implementation details, benchmark construction, and evaluation evidence. Section~\ref{sec:conclusion} concludes with limitations and future work.
